% ============ 第2章 §2.14 极值问题(原书 pp.158-161) ============
\section{极值问题}\label{sec:2.14}
\index{极值问题}
\index{极值函数}
\index{极值核}
\index{共轭关系}


当 $\phi \in (H^{p})^{*}$ 时,对
\begin{equation*}
\|\phi\| = \sup \{|\phi(f)|; f \in H^{p}, \|f\|_{p} \leq 1\}
\end{equation*}
的极值问题指的是下面的 (i)--(iv).(i) 上确界的值是什么?(ii) 上确界能否成为最大值?(iii) 给出最大值的函数(极值函数)是否唯一?(iv) 极值函数是什么样的函数?

一般的 $\phi \in (H^{p})^{*}$($1 \leq p < \infty$)可表示为
\begin{align*}
\phi(f) & = \frac{1}{2\pi}\int_{-\pi}^{\pi} f(e^{it}) e^{it} k(e^{it})\dd t \\
& = \frac{1}{2\pi i}\int_{|z|=1} f(z) k(z)\dd z \quad (f \in H^{p}),
\end{align*}
其中 $k \in L^{q}$ 且 $1/p + 1/q = 1$.$\phi \in (H^{\infty})^{*}$ 一般不能用 $k \in L^{1}$ 表示.这是因为 $(L^{\infty})^{*} \supsetneq L^{1}$.但显然用 $k \in L^{1}$ 表示的 $\phi$ 属于 $(H^{\infty})^{*}$.如 2.13 节所示,当 $1 \leq p < \infty$ 时 $(H^{p})^{*} = L^{q}/zH^{q}$,当 $p = \infty$ 时 $(H^{\infty})^{*} = (L^{1}/zH^{1})^{**} \supset L^{1}/zH^{1}$,因此若 $h - k \in H^{q}$,则 $h$ 与 $k$ 在 $(H^{p})^{*}$ 中给出同一元素.于是 $\|\phi\| \leq \|k + g\|_{q}$($g \in H^{q}$),但可以证明下面的共轭关系(定理~\ref{thm:2.14.2}),进而可以解 $1 < p < \infty$ 情形的极值问题 (ii).

\begin{lemma}\label{lem:2.14.1}
设 $X$ 为 Banach 空间,$S$ 为 $X$ 的闭子空间且 $S^{\perp} = \{\phi \in X^{*}; \phi(x) = 0 \ (x \in S)\}$.

(1) 若 $\phi \in X^{*}$,则
\begin{equation*}
\sup \{|\phi(x)|; x \in S, \|x\| \leq 1\} = \min \{\|\phi + \psi\|; \psi \in S^{\perp}\}.
\end{equation*}

(2) 若 $x \in X$,则
\begin{equation*}
\max \{|\psi(x)|; \psi \in S^{\perp}, \|\psi\| \leq 1\} = \inf \{\|x + y\|; y \in S\}.
\end{equation*}
\end{lemma}

\begin{proof}
(1) 定理~\ref{thm:2.13.1} 的证明是通用的,$S^{*} = X^{*}/S^{\perp}$ 成立.其证明即证明了 (1).(2) 定理~\ref{thm:2.13.2} 的证明是通用的,$S^{\perp} = (X/S)^{*}$ 成立.设 $\psi \in S^{\perp}$ 且 $\|\psi\| \leq 1$,则对任意 $x \in X$ 和任意 $y \in S$,
\begin{equation*}
|\psi(x)| = |\psi(x + y)| \leq \|x + y\|,
\end{equation*}
所以
\begin{equation*}
\sup \{|\psi(x)|; \psi \in S^{\perp}, \|\psi\| \leq 1\} \leq \|x + S\|.
\end{equation*}
由 Hahn--Banach 定理的推论 [4,引理 30.1],存在 $\Psi_{0} \in (X/S)^{*}$,使得 $\|\Psi_{0}(x+S)\| = \|x+S\|$ 且 $\|\Psi_{0}\| = 1$.令 $\psi_{0}(x) = \Psi_{0}(x+S)$,则 $\psi_{0} \in S^{\perp}$ 且 $\|\psi_{0}\| = \|\Psi_{0}\| = 1$,故 (2) 得证.
\end{proof}

\begin{theorem}\label{thm:2.14.2}
设 $1 \leq p \leq \infty$ 且 $1/p + 1/q = 1$.若 $k \in L^{q}$,则
\begin{align*}
\sup \left\{\left| \frac{1}{2\pi i}\int_{|z|=1} f(z) k(z)\dd z \right|; f \in H^{p}, \|f\|_{p} \leq 1 \right\}&= \min \{\|k - g\|_{q}; g \in H^{q}\}.
 \end{align*}
$p \neq 1$ 时 sup 可改为 max.
\end{theorem}

\begin{proof}
当 $1 < p \leq \infty$ 时,取 $X = L^{q}$,$S^{\perp} = zH^{p}$ 且 $S = H^{q}$,用引理~\ref{lem:2.14.1} 的 (2) 得
\begin{align*}
\max \left\{\left| \frac{1}{2\pi i}\int_{|z|=1} f(z) k(z)\dd z \right|; f \in H^{p}, \|f\|_{p} \leq 1 \right\}&= \inf \{\|k - g\|_{q}; g \in H^{q}\}.
 \end{align*}
另外,取 $X = L^{p}$,$S^{\perp} = H^{q}$ 且 $S = zH^{p}$,用引理~\ref{lem:2.14.1} 的 (1),inf 可改为 min.当 $p = 1$ 时,取 $X = L^{1}$,$S = zH^{1}$ 且 $S^{\perp} = H^{\infty}$,用引理~\ref{lem:2.14.1} 的 (1) 即可.
\end{proof}

共轭关系中给出最小值的函数称为极值核.

\begin{lemma}\label{lem:2.14.3}
设 $1 \leq p \leq \infty$ 且 $1/p + 1/q = 1$.$F \in H^{p}$ 为 $\phi \in (H^{p})^{*}$ 的极值函数且 $\phi(F) > 0$,其共轭关系的极值核为 $K$ 的充要条件如下:
\begin{equation*}
e^{i\theta} F(e^{i\theta}) K(e^{i\theta}) \geq 0 \quad \text{a.e. } \theta,
\end{equation*}
并且,若 $p = 1$,则 a.e. $\theta$ 有 $|K(e^{i\theta})| = \|K\|_{\infty}$;若 $1 < p < \infty$,则 a.e. $\theta$ 有 $|F(e^{i\theta})|^{p} = \|K\|_{q}^{-q}|K(e^{i\theta})|^{q}$;若 $p = \infty$,则在 $\{\theta; K(e^{i\theta}) \neq 0\}$ 上 a.e. $\theta$ 有 $|F(e^{i\theta})| = 1$.
\end{lemma}

\begin{proof}
由定理~\ref{thm:2.14.2},
\begin{equation*}
\frac{1}{2\pi}\int_{-\pi}^{\pi} F(e^{i\theta}) e^{i\theta} K(e^{i\theta})\dd \theta = \|K\|_{q}\|F\|_{p}.
\end{equation*}
由 Hölder 不等式,
\begin{equation*}
\frac{1}{2\pi}\int_{-\pi}^{\pi} e^{i\theta} F(e^{i\theta}) K(e^{i\theta})\dd \theta = \frac{1}{2\pi}\int_{-\pi}^{\pi} |e^{i\theta} F(e^{i\theta}) K(e^{i\theta})|\dd \theta,
\end{equation*}
故 $e^{i\theta}F(e^{i\theta})K(e^{i\theta}) \geq 0$ a.e. $\theta$.又因 Hölder 不等式中等式成立,后半部分关于 $|K|$ 与 $|F|$ 的关系式得证.
\end{proof}

\begin{theorem}\label{thm:2.14.4}
设 $1 \leq p \leq \infty$ 且 $\phi \in (H^{p})^{*}$.

(1) 当 $1 < p < \infty$ 时,$H^{p}$ 的极值问题的解存在,且除常数倍外唯一.

(2) 当 $p = \infty$ 时,若 $\phi \in L^{1}/zH^{1}$,则 $H^{\infty}$ 的极值问题的解存在,且除常数倍外唯一.

(3) 当 p = 1 时,若 $\phi \in C/zA$,则 $H^{1}$ 的极值问题的解存在,但不一定除常数倍外唯一.
\end{theorem}

\begin{proof}
(1) 与 (2) 中,解的存在性是定理~\ref{thm:2.14.2} 的结果,解的唯一性由引理~\ref{lem:2.14.3} 得出.(3) 中,若 $\phi \in C/zA$,则 $\phi$ 是 *弱连续的,$H^{1}$ 的单位球是 *弱紧的,故解存在.
\end{proof}

若
\begin{equation*}
\phi(f) = \frac{1}{2\pi i}\int_{|z|=1} f(z)\bar{z}^{2}\dd z,
\end{equation*}
则解有无穷多个.
